\section{Multilevel headers and spanning titles}
\key{\gradheader{...}} adds a header layer. Within it,
\key{\gradspan{count}{title}} spans consecutive columns. Declare all columns
first, then header layers, then groups and rows. Each layer must span exactly
the total number of columns. Column titles form the bottom header layer.
\begin{demo}{Three header levels}
\begin{gradtable}[rules=grid,header align=c,
  header color=gradbarsBlue!8,text width=28mm,
  column sep=8pt,row padding=3pt]
  \gradcolumn{method}{Method}
  \gradcolumn[type=number,width=24mm,
    options={precision=1}]{accuracy}{Accuracy}
  \gradcolumn[type=number,width=24mm,
    options={precision=1}]{fscore}{F1}
  \gradcolumn[type=number,width=24mm,
    options={precision=0}]{latency}{ms}
  \gradheader{\gradspan{4}{Model evaluation}}
  \gradheader{\gradspan{1}{Setup}\gradspan{2}{Prediction quality}
    \gradspan{1}{Latency}}
  \gradrow{A,91.4,89.7,24}
  \gradrow{B,93.2,91.5,31}
\end{gradtable}
\end{demo}
Spanning titles are centered, follow header color and use the chosen rules.
An empty title still occupies its declared span. Layers merge horizontally,
not vertically. Schemas may include headers, but no columns may follow a header.
